\documentclass[11pt]{article}
\usepackage[margin=1in]{geometry}
\usepackage{amsmath,amssymb}
\usepackage{booktabs}
\usepackage{enumitem}
\usepackage[hidelinks]{hyperref}
\usepackage{xcolor}
\hypersetup{colorlinks=true, linkcolor=blue!60!black, urlcolor=blue!60!black}

\newcommand{\eb}{\bar{e}}
\newcommand{\SJ}{\mathit{SJ}}
\newcommand{\ReHQJ}{\mathit{ReHQJ}}

\title{ReHQJ: the HexOpt Paper and conspire}
\author{autotwin/automesh issue \#768}
\date{October 14, 2026}

\begin{document}
\maketitle

\section*{Sources}

This note compares the HexOpt objective in two sources.

\begin{itemize}[leftmargin=*]\sloppy
  \item \textbf{Paper.} H.~Tong and Y.~J.~Zhang, ``Fast and Robust Hexahedral Mesh Optimization via Augmented Lagrangian, L-BFGS, and Line Search,'' \href{https://arxiv.org/pdf/2410.11656}{arXiv:2410.11656}, Equations~(2.5), (2.6), (2.8), and (2.9).
  \item \textbf{conspire.} Branch \texttt{msj} at commit \texttt{515c78cee}, draft \href{https://github.com/mrbuche/conspire.rs/pull/251}{mrbuche/conspire.rs\#251}, file \texttt{src/\allowbreak geometry/\allowbreak mesh/\allowbreak quality/\allowbreak improve/\allowbreak msj/\allowbreak mod.rs}, functions \texttt{evaluate} and \texttt{terms}.
\end{itemize}

$h$ is one hex element of the mesh $H$, and $\eb$ is its average edge length.
$\SJ(h)$ is the minimum of the scaled Jacobian over nine points of $h$: its eight corners and its body center (paper, Section~2.1).
$J(h)$ is the minimum of the Jacobian over the same nine points.
conspire computes both minimums in \texttt{terms}, each over all nine points.
$\Theta$ is the threshold.
Both sources hold $\eb$ fixed when they take the gradient.

\section*{Definitions}

The paper calls the per-hex term $\ReHQJ(h, \Theta)$.
This note writes it as $\psi(h)$, with subscript $\mathrm{p}$ for the paper and $\mathrm{c}$ for conspire.
Both sources use the same term for an inverted hex:
\begin{equation*}
  \psi_{\mathrm{p}}(h) = \psi_{\mathrm{c}}(h) = J(h) / \eb, \qquad J(h) \le 0.
\end{equation*}
For $J(h) > 0$, the paper defines
\begin{equation*}
  \psi_{\mathrm{p}}(h) =
  \begin{cases}
    \eb^2 \, \SJ(h), & \SJ(h) \le \Theta, \\
    \Theta, & \SJ(h) > \Theta,
  \end{cases}
\end{equation*}
and conspire implements
\begin{equation*}
  \psi_{\mathrm{c}}(h) = \eb^2 \min\bigl(\SJ(h), \Theta\bigr).
\end{equation*}
The two differ only above $\Theta$, where $\psi_{\mathrm{p}} = \Theta$ and $\psi_{\mathrm{c}} = \Theta \, \eb^2$.
$\psi_{\mathrm{c}}$ is one expression: $\eb^2$ times $\SJ$ clamped at $\Theta$.
$\psi_{\mathrm{p}}$ has no such form unless $\eb = 1$.

\section*{Units}

$J = \det(e_0, e_1, e_2)$, where each edge vector has units of length.
$\SJ$ and $\Theta$ have no units.

\begin{center}
\begin{tabular}{lcc}
  \toprule
  case & $\psi_{\mathrm{p}}$ & $\psi_{\mathrm{c}}$ \\
  \midrule
  $J \le 0$ & length$^2$ & length$^2$ \\
  $J > 0$, $\SJ \le \Theta$ & length$^2$ & length$^2$ \\
  $\SJ > \Theta$ & \textbf{none} & length$^2$ \\
  \bottomrule
\end{tabular}
\end{center}

\section*{The minimized function}

At each threshold $\Theta$, conspire minimizes
\begin{equation*}
  L = -\sum_{h \in H} \psi_{\mathrm{c}}(h) + \sum_{i} \Bigl[ \lambda_i \, r_i + \frac{\rho}{2} \, r_i^2 \Bigr].
\end{equation*}
\begin{itemize}[leftmargin=*]
  \item The first sum has one term per hex, with $\psi_{\mathrm{c}}$ as defined above.
  \item The second sum is an augmented Lagrangian that holds the boundary on the target STL.
    $r_i = n_i \cdot x - d_i$ is the signed distance from a boundary node $x$ to the plane of one STL triangle.
    Each boundary node gets up to four planes: the plane of its closest triangle, plus the planes of nearby triangles whose normals differ from those already kept by more than $30^\circ$.
  \item After each L-BFGS run, conspire updates each multiplier, $\lambda_i \leftarrow \lambda_i + \rho \, r_i$.
    $\rho$ starts at 100 and doubles each time the mesh meets $\Theta$ with the boundary on the STL.
  \item $r_i$ and $\lambda_i$ have units of length, and $\rho$ has none.
    Both constraint terms have units of length$^2$.
\end{itemize}
The paper's Equation~(2.8) has the same structure.
Its constraint is $Z_k - Z^t_k$, the offset from each boundary vertex to its closest point on the target surface.
Its Equation~(2.9) updates the multipliers the same way, and it doubles $\rho$ each time the minimum scaled Jacobian reaches $\Theta$.
The paper's $L$ uses $\psi_{\mathrm{p}}$ in place of $\psi_{\mathrm{c}}$. The next section returns to that difference.

\section*{Behavior at the threshold}

\begin{enumerate}[leftmargin=*]
  \item \textbf{Continuity.}
    At $\SJ = \Theta$, both forms below $\Theta$ give $\Theta \, \eb^2$.
    Above $\Theta$, $\psi_{\mathrm{c}}$ gives the same value, so it is continuous at $\Theta$.
    Above $\Theta$, $\psi_{\mathrm{p}}$ gives $\Theta$, so it jumps by $\Theta (\eb^2 - 1)$ unless $\eb = 1$.
    Neither form is differentiable at $\Theta$.
    The slope drops to zero there in both, by design.

  \item \textbf{The sign of the jump.}
    $L$, above, sums $-\psi$ over all hexes.
    Under $\psi_{\mathrm{p}}$, a hex that rises above $\Theta$ changes $L$ by $+\Theta (\eb^2 - 1)$.
    For $\eb > 1$, $L$ rises.
    The Armijo condition requires $L$ to fall, so the jump counts against every step that lifts such a hex above $\Theta$.
    For $\eb < 1$, $L$ falls, so the jump rewards the crossing instead.
    The sign depends on the length unit of the mesh.
    \begin{itemize}
      \item On the snapped bone octree in the reference frame, $\eb$ runs from 0.644 to 2.81, with a median of 1.253.
        The same mesh has hexes on both sides of $\eb = 1$.
      \item On the scale-5 bone mesh in the STL's own units, $\eb$ runs from 0.009 to 0.034.
    \end{itemize}
    Under $\psi_{\mathrm{c}}$, the crossing changes $L$ by zero.

  \item \textbf{Units.}
    $\psi_{\mathrm{c}}$ has units of length$^2$ in every case.
    conspire's types enforce this: each term is a \texttt{Quantity<Area>}.
\end{enumerate}

For a hex that stays above $\Theta$ through a step, both forms are constants.
They give the same gradient and the same energy differences.
The forms differ only when a hex crosses $\Theta$.

\end{document}
